\chapter{Validation and Acceptance}
\label{ch:validation}

\section{Rule-pack validation}

Every rule-pack pair must be backed by an executable validation harness.
The harness must prove at least:
\begin{itemize}
  \item the files exist at the intended paths;
  \item required sections are present;
  \item referenced source files and schema files exist;
  \item threshold values and condition codes are explicit;
  \item the pair benchmarks at least as strongly as the established
    operational baseline.
\end{itemize}

\section{Recommended benchmark baseline}

The current repository baseline for operational specificity is the
Generative UI deployment rule pack at
\filepath{src/generativeUI/.clinerules}. Any new domain pack should be
benchmarked against that file for:
\begin{itemize}
  \item incident memory;
  \item concrete verification steps;
  \item threshold explicitness;
  \item checklist quality;
  \item operational specificity.
\end{itemize}

\section{Quantified benchmark scorecard}
\label{sec:benchmark-scorecard}

To enable objective, reproducible comparison against the baseline, the
validation harness SHOULD compute a numeric score based on the following
metrics. This scorecard allows new domain packs to self-assess whether
they meet the ``at least as strong'' requirement.

\subsection{Scorecard metrics}

\begin{table}[htbp]
\centering
\footnotesize
\caption{Operational specificity scorecard.}
\label{tab:scorecard}
\begin{tabularx}{\textwidth}{@{}L{4.5cm} l X@{}}
\toprule
\textbf{Metric} & \textbf{Pass threshold} & \textbf{Calculation method} \\
\midrule
Explicit numeric thresholds & 100\% & All monitoring rules must specify numeric values (e.g., ``> 5\%''), not vague terms (e.g., ``high rate''). \\
Verification path density & $\ge 5$ paths per 100 lines & Count of explicit file paths referenced divided by rule pack line count, normalized to 100 lines. \emph{Note: Shorter files (< 60 lines) must contain at least 3 unique paths to avoid penalty.} \\
Known issue registry entries & $\ge 3$ entries & Count of documented previous incidents, workarounds, or known edge cases. \\
Checklist completeness & 100\% & All 11 development pack sections (or 9 runtime pack sections) must be present and non-empty. \\
Schema reference coverage & 100\% & All payload types mentioned must reference a schema file in \filepath{docs/schema/}. \\
Test entry point count & $\ge 2$ per pack & Count of explicit test file paths or commands referenced. \\
\bottomrule
\end{tabularx}
\end{table}

\subsection{Scoring algorithm}

The validation harness SHALL compute scores as follows:

\begin{enumerate}
  \item \textbf{Parse the rule pack} into sections using header markers.
  \item \textbf{Extract metrics}:
    \begin{itemize}
      \item Count threshold expressions matching pattern
        \verb|[><=]+\s*\d+(\.\d+)?[%]?|.
      \item Count file path references matching pattern
        \verb|src/|, \verb|docs/|, or \verb|tests/|.
      \item Count known issue entries in the ``Known Issues'' or
        ``Current Repo Reality'' section.
      \item Count schema references matching pattern
        \verb|\.schema\.(json|yaml)|.
      \item Count test references matching pattern
        \verb|test_.*\.py| or \verb|make test|.
    \end{itemize}
  \item \textbf{Compare against baseline}: Read the
    \filepath{.baseline-clinerules-pointer} file in the repository root
    to identify the current gold-standard pack (e.g.,
    \filepath{src/generativeUI/.clinerules}). Load that pack and compute
    the same metrics.
  \item \textbf{Pass/fail determination}: A pack passes if all metrics
    meet or exceed the thresholds in \cref{tab:scorecard} AND are not
    worse than the baseline values.
\end{enumerate}

\section{The .baseline-clinerules-pointer}

To ensure the benchmark evolves as engineering standards mature, the
validation harness SHALL NOT hardcode the baseline pack path. This
pointer file allows the platform team to rotate the baseline without
requiring a specification update or code change to the harness.

\subsection{Scorecard output format}

The harness SHALL produce a JSON report:

\begin{verbatim}
{
  "pack_path": "src/training/.clinerules",
  "baseline_path": "src/generativeUI/.clinerules",
  "timestamp": "2026-04-21T09:00:00Z",
  "metrics": {
    "threshold_explicitness": {"value": 100, "pass": true},
    "path_density": {"value": 7.2, "pass": true},
    "known_issues": {"value": 5, "pass": true},
    "checklist_completeness": {"value": 100, "pass": true},
    "schema_coverage": {"value": 100, "pass": true},
    "test_entry_points": {"value": 3, "pass": true}
  },
  "overall_pass": true,
  "baseline_comparison": "EQUAL_OR_STRONGER"
}
\end{verbatim}

\subsection{Scorecard integration}

\begin{itemize}
  \item The scorecard MUST be computed as part of the CI/CD validation
    gate for any PR that modifies a \filepath{.clinerules} file.
  \item Scorecard results MUST be attached to the PR as a comment or
    status check.
  \item A pack that fails the scorecard MUST NOT be merged without
    explicit exception approval (see \cref{sec:conflict-detection}).
\end{itemize}

\section{Acceptance commands}

\begin{table}[htbp]
\centering
\footnotesize
\caption{Minimum acceptance commands.}
\label{tab:acceptance}
\begin{tabularx}{\textwidth}{@{}l X@{}}
\toprule
\textbf{Area} & \textbf{Command or gate} \\
\midrule
Training rule packs & \texttt{cd src/training \&\& make test-clinerules} \\
Regulations rule packs & \texttt{cd src/intelligence/ai-core-pal \&\& make test-clinerules} \\
Regulations MCP surface & \texttt{cd src/intelligence/ai-core-pal \&\& make test} \\
Agent supplement PDF & \texttt{make spec-clinerules-agents} \\
PDF publication & \texttt{make spec-<name>} for each published book \\
\bottomrule
\end{tabularx}
\end{table}

\section{Publication criteria}

The swarm may claim a spec is delivered only when:
\begin{enumerate}
  \item the relevant rule packs exist;
  \item the harnesses pass;
  \item the LaTeX source is committed and reference-clean;
  \item the build target exists for the corresponding PDF output;
  \item the runtime-monitor contract is present for any domain that is
    expected to run live.
\end{enumerate}

\section{Failure conditions}

\begin{itemize}
  \item Missing runtime-monitor pack for a governed live domain.
  \item Rule pack cites non-existent source of truth or non-existent
    verification paths.
  \item Domain pack conflicts with regulations obligations.
  \item PDF build succeeds but the supporting rule-pack harness fails.
\end{itemize}
